\section{延伸问题：非线性规划}

机械臂运行速度提高时，单个任务的搬运时间会缩短，系统最大完工时间 $C_{\max}$ 可能下降；但较高速度通常也会带来更大的驱动能耗、协调能耗和控制代价。相反，若机械臂运行速度降低，系统能耗可能下降，但任务持续时间会变长，甚至可能无法满足给定的截止时间要求。

这种处理方式可以理解为\textbf{先离散调度，后连续调速}。前一层由整数规划、目标规划和隐枚举等方法确定调度结构，后一层在调度结构已知的条件下，通过非线性规划优化速度参数。\textbf{既避免了将离散变量和连续变量完全耦合所带来的求解困难，又能够进一步分析速度、完成时间和能耗之间的权衡关系。}


主问题已经给出了某一系统模式下的基准调度结果，仅将速度倍率作为连续决策变量。

设单臂任务速度倍率为 $s_1$，双臂协同任务速度倍率为 $s_2$，二者满足速度上下界约束：
\[
	s_{\min} \leq s_1 \leq s_{\max},
	\qquad
	s_{\min} \leq s_2 \leq s_{\max}.
\]
其中，$s_1=1$ 和 $s_2=1$ 表示采用主问题中的基准速度；$s_1>1$ 或 $s_2>1$ 表示提高对应任务类型的运行速度；$s_1<1$ 或 $s_2<1$ 表示降低对应任务类型的运行速度。

在调度结构已知的情况下，系统最大完工时间可以表示为速度倍率的函数。由于任务执行时间与速度近似成反比，因此可写为：
\[
	C_{\max}(s_1,s_2)
	=
	C_0
	+
	\frac{T_1}{s_1}
	+
	\frac{T_2}{s_2}.
\]
其中，$C_0$ 表示与速度倍率无关的固定时间部分，$T_1$ 表示单臂任务对应的可压缩时间部分，$T_2$ 表示双臂协同任务对应的可压缩时间部分。

能耗函数则表示为速度倍率的非线性函数：
\[
	E(s_1,s_2)
	=
	E_0
	+
	a_1 s_1^{\alpha}
	+
	a_2 s_2^{\beta}
	+
	\rho_1(s_1-1)^2
	+
	\rho_2(s_2-1)^2 .
\]
其中，$E_0$ 表示与速度倍率无关的固定能耗部分，$a_1 s_1^{\alpha}$ 和 $a_2 s_2^{\beta}$ 分别表示单臂任务和双臂协同任务的速度相关能耗，$\alpha$ 和 $\beta$ 为非线性指数，$\rho_1(s_1-1)^2$ 和 $\rho_2(s_2-1)^2$ 表示速度偏离基准速度时带来的附加代价。

在给定截止时间 $D$ 的条件下，两个延伸问题都需要满足：
\[
	C_{\max}(s_1,s_2) \leq D.
\]
因此，共同的非线性优化思想可以概括为：
\[
	\min_{s_1,s_2} \quad E(s_1,s_2)
\]
\[
	\mathrm{s.t.} \quad
	C_{\max}(s_1,s_2) \leq D,
\]
\[
	s_{\min} \leq s_1 \leq s_{\max},
	\qquad
	s_{\min} \leq s_2 \leq s_{\max}.
\]

上述模型说明，本章非线性规划部分的核心不再是选择哪只机械臂执行任务，而是在调度方案已经确定之后，进一步\textbf{选择合适的运行速度}。后续两个问题将在这一共同模型基础上分别进行扩展。



\subsection{问题一：速度--能耗非线性优化问题}

\subsubsection{问题分析}

不同 seed 对应不同的物块空间分布，因此会导致搬运距离、任务衔接关系、最大完工时间和系统能耗产生一定波动。若直接选取某一个 seed 进行速度优化，结果可能带有较强的偶然性，不能代表该任务规模下的整体表现。因此首先对多个 seed 的主问题结果进行平均，得到二臂和三臂模式下的平均基准完工时间与平均基准能耗。

设某一模式下的平均基准最大完工时间为 $C_{\max}^{0}$，平均基准能耗为 $E^{0}$。根据任务类型数量将基准时间和基准能耗拆分为单臂任务部分和双臂协同任务部分。

设 Type1 和 Type2 的任务数量分别为 $n_1,n_2$，Type3 和 Type4 的任务数量分别为 $n_3,n_4$。单臂任务占比和双臂任务占比分别定义为
\[
	\omega_1
	=
	\frac{n_1+n_2}{n_1+n_2+n_3+n_4},
	\qquad
	\omega_2
	=
	\frac{n_3+n_4}{n_1+n_2+n_3+n_4}.
\]
于是基准完工时间可近似拆分为
\[
	T_1=\omega_1 C_{\max}^{0},
	\qquad
	T_2=\omega_2 C_{\max}^{0}.
\]
其中，$T_1$ 表示单臂任务对应的可压缩时间部分，$T_2$ 表示双臂协同任务对应的可压缩时间部分。

同理，基准能耗可拆分为
\[
	E_1^{0}=\omega_1 E^{0},
	\qquad
	E_2^{0}=\omega_2 E^{0}.
\]
其中，$E_1^{0}$ 表示单臂任务对应的基准能耗，$E_2^{0}$ 表示双臂协同任务对应的基准能耗。

在此基础上，引入单臂任务速度倍率 $s_1$ 和双臂协同任务速度倍率 $s_2$。当速度倍率改变后，最大完工时间表示为
\[
	C_{\max}(s_1,s_2)
	=
	\frac{T_1}{s_1}
	+
	\frac{T_2}{s_2}.
\]
该式表示速度倍率越大，对应任务类型的执行时间越短。

能耗模型：
\[
	E(s_1,s_2)
	=
	E_1^{0}
	\left[
	\frac{\lambda_1}{s_1}
	+
	(1-\lambda_1)
	\left((1-\rho_1)+\rho_1 s_1^2\right)
	\right]
	+
	E_2^{0}
	\left[
	\frac{\lambda_2}{s_2}
	+
	(1-\lambda_2)
	\left((1-\rho_2)+\rho_2 s_2^2\right)
	\right].
\]
其中，$\lambda_1,\lambda_2$ 用于刻画低速运行时因任务持续时间增加而产生的保持能耗影响，$\rho_1,\rho_2$ 用于刻画高速运行时驱动能耗增加的影响。该模型并不是简单认为速度越低能耗越小，而是认为存在一个相对经济的速度区间：速度过低会拉长执行时间，速度过高会增加驱动代价。

对于给定截止时间 $D$，问题一最终可写为
\[
	\min_{s_1,s_2} \quad E(s_1,s_2)
\]
\[
	\mathrm{s.t.} \quad
	C_{\max}(s_1,s_2) \leq D,
\]
\[
	s_{\min} \leq s_1 \leq s_{\max},
	\qquad
	s_{\min} \leq s_2 \leq s_{\max}.
\]

其中，截止时间 $D$ 由基准完工时间乘以不同截止时间比例得到。通过改变截止时间比例，可以观察在宽松工期和紧张工期下，二臂模式与三臂模式的最优速度、最小能耗和推荐结果如何变化。因此，问题一的核心不是重新安排任务，而是在平均调度结果已经确定的条件下，进一步回答“在不同工期要求下，应该以多快的速度运行才更节能”。



\subsubsection{实现过程}

本问题属于带不等式约束的连续非线性规划问题，决策变量为单臂任务速度倍率 $s_1$ 和双臂协同任务速度倍率 $s_2$。由于速度倍率必须为正，且存在上下界约束和截止时间约束，因此本文采用内点法处理约束，并在每次确定下降方向后使用 0.618 法进行一维步长搜索。

首先说明该问题满足使用上述方法的基本条件。根据前文建模，系统完工时间函数可写为
\[
	C_{\max}(s_1,s_2)
	=
	\frac{T_1}{s_1}
	+
	\frac{T_2}{s_2}.
\]
在 $s_1>0,\ s_2>0$ 的定义域内，函数 $1/s$ 是凸函数，因此 $C_{\max}(s_1,s_2)$ 也是凸函数。截止时间约束
\[
	C_{\max}(s_1,s_2)\leq D
\]
对应凸函数的下水平集，因此形成凸可行域。

能耗函数采用 U 形经济速度模型。对任一任务类型，其速度相关能耗项可写为
\[
	f_i(s)
	=
	\frac{\lambda_i}{s}
	+
	(1-\lambda_i)\left[(1-\rho_i)+\rho_i s^2\right],
	\qquad i=1,2.
\]
对其求二阶导数，有
\[
	f_i''(s)
	=
	\frac{2\lambda_i}{s^3}
	+
	2(1-\lambda_i)\rho_i .
\]
由于 $s>0$，且 $\lambda_i,\rho_i$ 均为非负参数，因此
\[
	f_i''(s)\geq 0.
\]
这说明单个速度能耗项为凸函数。总能耗函数由若干凸函数和常数项加权相加得到，因此仍为凸函数。由此，本问题可以看作在凸可行域上最小化凸目标函数的非线性规划问题。该性质保证了局部最优解具有较好的全局解释，也使得内点法和一维搜索具有合理的使用基础。

为使用内点法，将所有不等式约束统一写成 $g_i(s_1,s_2)>0$ 的形式。具体定义为
\[
	g_0(s_1,s_2)
	=
	D-C_{\max}(s_1,s_2),
\]
\[
	g_1(s_1,s_2)=s_1-s_{\min},
	\qquad
	g_2(s_1,s_2)=s_{\max}-s_1,
\]
\[
	g_3(s_1,s_2)=s_2-s_{\min},
	\qquad
	g_4(s_1,s_2)=s_{\max}-s_2.
\]
其中，$g_0$ 对应截止时间约束，$g_1$ 到 $g_4$ 对应两个速度变量的上下界约束。只要所有 $g_i>0$，当前点就严格位于可行域内部。

内点法的核心思想是将原来的约束优化问题转化为无约束优化问题。本文构造如下对数障碍函数：
\[
	\Phi_{\mu}(s_1,s_2)
	=
	E(s_1,s_2)
	-
	\mu
	\sum_{i=0}^{4}
	\ln g_i(s_1,s_2),
\]
其中，$\mu>0$ 为障碍参数。由于 $\ln g_i$ 只有在 $g_i>0$ 时才有定义，因此迭代点会被限制在可行域内部。当某个约束即将被违反时，对应的 $g_i$ 接近 0，$-\ln g_i$ 会迅速增大，从而阻止迭代点越过可行域边界。随着 $\mu$ 逐渐减小，障碍项影响减弱，障碍函数的最优点逐渐逼近原约束问题的最优解。

在每一次内点迭代中，程序先根据当前速度点
\[
	\boldsymbol{s}^{(k)}
	=
	\begin{bmatrix}
		s_1^{(k)}\
		s_2^{(k)}
	\end{bmatrix}
\]
计算障碍函数的梯度 $\nabla \Phi_{\mu}(\boldsymbol{s}^{(k)})$，并取负梯度方向作为下降方向：
\[
	\boldsymbol{d}^{(k)}
	=
	-
	\nabla \Phi_{\mu}(\boldsymbol{s}^{(k)}).
\]
这样，原来的二维搜索问题就转化为沿着方向 $\boldsymbol{d}^{(k)}$ 寻找合适步长的问题。新的速度点表示为
\[
	\boldsymbol{s}^{(k+1)}
	=
	\boldsymbol{s}^{(k)}
	+
	\alpha \boldsymbol{d}^{(k)},
\]
其中 $\alpha$ 是待确定的一维步长。

0.618 法本身是一维搜索方法，不能直接用于二维变量 $(s_1,s_2)$。\textbf{代码先由内点法确定二维空间中的下降方向 $\boldsymbol{d}^{(k)}$，再把该方向上的步长 $\alpha$ 转化为一维搜索问题。}于是构造一维函数
\[
	\varphi(\alpha)
	=
	\Phi_{\mu}
	\left(
	\boldsymbol{s}^{(k)}
	+
	\alpha \boldsymbol{d}^{(k)}
	\right).
\]
此时，$\varphi(\alpha)$ 只含一个变量 $\alpha$，可以使用 0.618 法进行求解。

由于障碍函数 $\Phi_{\mu}$ 在可行域内部保持凸性，而凸函数沿任意直线方向的限制函数仍为一维凸函数，因此 $\varphi(\alpha)$ 在可行步长区间内具有单峰性。0.618 法正适用于这种一维单峰函数的最小值搜索。\textbf{这里不是直接用 0.618 法求二维最优点，而是采用“内点法定方向 + 0.618 法定步长”的组合结构。}

在实际实现中，首先根据速度上下界和截止时间约束确定可行步长区间
\[
	0\leq \alpha \leq \alpha_{\max},
\]
并保证对任意候选步长都有
\[
	g_i
	\left(
	\boldsymbol{s}^{(k)}
	+
	\alpha \boldsymbol{d}^{(k)}
	\right)
	>
	0,
	\qquad i=0,1,\ldots,4.
\]
然后在该区间内使用 0.618 法不断缩小搜索范围，比较两个试探点处的 $\varphi(\alpha)$ 值，保留函数值较小的一侧区间，直到区间长度小于给定精度为止。最终得到当前方向上的近似最优步长 $\alpha^{(k)}$。

这个过程如图 \ref{fig:interior-golden-flow} 所示。

\begin{figure}[htbp]
	\centering
	\begin{tikzpicture}[
		node distance=0.45cm,
		flowbox/.style={
			rectangle,
			rounded corners=2pt,
			draw=gray!55,
			line width=0.4pt,
			minimum height=0.72cm,
			text width=13.5cm,
			align=left,
			inner xsep=8pt,
			inner ysep=5pt,
			font=\small
		},
		arrow/.style={
			-{Stealth[length=2mm,width=1.5mm]},
			draw=gray!70,
			line width=0.4pt
		}
		]
		
		\node[flowbox, fill=purple!8] (step1)
		{\textbf{读取基准数据}\quad 读取平均基准完工时间、平均基准能耗和任务数量组合，构造 $C_{\max}(s_1,s_2)$ 与 $E(s_1,s_2)$};
		
		\node[flowbox, fill=green!8, below=of step1] (step2)
		{\textbf{选取初始内点}\quad 给定截止时间 $D$ 和速度上下界 $s_{\min},s_{\max}$，选取满足所有约束的初始内点};
		
		\node[flowbox, fill=purple!8, below=of step2] (step3)
		{\textbf{构造障碍函数}\quad 构造 $\Phi_{\mu}(s_1,s_2)$，将不等式约束转化为对数障碍项};
		
		\node[flowbox, fill=green!8, below=of step3] (step4)
		{\textbf{确定下降方向}\quad 在当前点计算 $\boldsymbol{d}^{(k)}=-\nabla \Phi_{\mu}(\boldsymbol{s}^{(k)})$};
		
		\node[flowbox, fill=purple!8, below=of step4] (step5)
		{\textbf{转化为一维搜索}\quad 将二维问题限制在下降方向上，得到一维函数 $\varphi(\alpha)$};
		
		\node[flowbox, fill=green!8, below=of step5] (step6)
		{\textbf{0.618 法搜索步长}\quad 在可行步长区间内搜索最优步长 $\alpha^{(k)}$};
		
		\node[flowbox, fill=purple!8, below=of step6] (step7)
		{\textbf{更新速度倍率}\quad 更新 $\boldsymbol{s}^{(k+1)}=\boldsymbol{s}^{(k)}+\alpha^{(k)}\boldsymbol{d}^{(k)}$};
		
		\node[flowbox, fill=green!8, below=of step7] (step8)
		{\textbf{迭代直到收敛}\quad 逐渐减小障碍参数 $\mu$，直到速度变化量、目标函数变化量或梯度范数满足停止条件};
		
		\draw[arrow] (step1) -- (step2);
		\draw[arrow] (step2) -- (step3);
		\draw[arrow] (step3) -- (step4);
		\draw[arrow] (step4) -- (step5);
		\draw[arrow] (step5) -- (step6);
		\draw[arrow] (step6) -- (step7);
		\draw[arrow] (step7) -- (step8);
		
	\end{tikzpicture}
	\caption{内点法结合 0.618 法求解速度--能耗非线性优化问题的流程}
	\label{fig:interior-golden-flow}
\end{figure}



通过上述方法，程序最终得到最优单臂速度倍率 $s_1^{*}$、最优双臂速度倍率 $s_2^{*}$、优化后的最大完工时间 $C_{\max}(s_1^{*},s_2^{*})$ 和优化后的总能耗 $E(s_1^{*},s_2^{*})$。该求解过程体现了内点法处理不等式约束、0.618 法处理一维步长搜索的结合：前者保证迭代始终在可行域内部进行，后者保证每一步沿下降方向取得较优步长，从而完成二维非线性规划问题的数值求解。


\subsubsection{KKT验证条件}

对于速度--能耗非线性优化问题，其标准形式可以写为
\[
\begin{aligned}
    \min_x\quad & f(x)\\
    \text{s.t.}\quad & g_j(x)\geq 0,\qquad j=1,2,\ldots,5.
\end{aligned}
\]
其中，$x=(s_S,s_D)^T$，$g_1$ 至 $g_4$ 分别对应速度上下界约束，$g_5$ 对应截止时间约束。

该问题的 KKT 条件为
\[
\nabla f(x^*)-\sum_{j=1}^{5}\mu_j \nabla g_j(x^*)=0,
\]
\[
g_j(x^*)\geq 0,\qquad j=1,2,\ldots,5,
\]
\[
\mu_j\geq 0,\qquad j=1,2,\ldots,5,
\]
\[
\mu_j g_j(x^*)=0,\qquad j=1,2,\ldots,5.
\]
其中，第一式为平稳性条件，第二式为原问题可行性条件，第三式为乘子非负条件，第四式为互补松弛条件。

本文对速度--能耗优化结果逐行检查上述条件。具体检查内容包括：各约束裕量 $g_1,g_2,g_3,g_4,g_5$ 是否满足非负约束；对应乘子 $\mu_1,\mu_2,\mu_3,\mu_4,\mu_5$ 是否满足非负条件；平稳性残差
\[
\left\|
\nabla f(x^*)-\sum_{j=1}^{5}\mu_j\nabla g_j(x^*)
\right\|
\]
是否足够小；以及最大互补误差
\[
\max_j \left|\mu_j g_j(x^*)\right|
\]
是否在数值容差范围内。


\subsection{问题二：鲁棒速度模式选择问题}

\subsubsection{问题分析}

问题二是在问题一基础上的进一步扩展。问题一对同一任务数量组合下的多个 seed 结果取平均，得到平均意义下的速度--能耗优化结果。这种处理方式能够反映某一任务规模的典型表现，但也会掩盖不同随机 seed 之间的差异。实际上，不同 seed 对应不同的物块空间分布，物块到收集盒的距离、任务之间的空载转移代价、中心安全区占用情况以及机械臂负载情况都会发生变化。因此，即使任务数量组合相同，不同 seed 下的基准完工时间和基准能耗也可能不同。

如果只根据平均结果选择速度倍率，可能会出现一种情况：该速度在平均意义下满足截止时间约束，但对于某些较困难的随机实例并不可行，或者在某些 seed 下能耗明显偏高。因此，问题二不再把多个 seed 简单平均，而是保留每个 seed 对应的独立调度结果，并要求同一任务数量组合下的所有 seed 共用同一组速度倍率。这样得到的速度方案不只适用于某一个随机实例，而是对同一任务规模下的多种空间分布都具有可行性。

该问题可以理解为鲁棒速度模式选择问题。这里的“鲁棒”主要体现在两个方面：第一，速度方案必须使所有 seed 都满足截止时间约束，而不是只满足平均情况；第二，优化目标不再是平均能耗最小，而是最坏情形能耗尽可能小。也就是说，模型关注的是在随机位置扰动下，二臂或三臂模式是否仍然能够稳定地满足工期要求，并控制最不利情况下的能耗水平。

设同一任务数量组合下的 seed 集合为 $\mathcal{R}$。对于任意 seed $r\in\mathcal{R}$，主问题已经分别给出了二臂或三臂模式下的基准调度结果。与问题一类似，本问题仍然不改变该 seed 下的任务分配和执行顺序，只对单臂任务速度倍率和双臂协同任务速度倍率进行优化。设单臂任务速度倍率为 $s_S$，双臂协同任务速度倍率为 $s_D$，其中下标 $S$ 表示 single-arm，下标 $D$ 表示 dual-arm。

对于每一个 seed $r$，根据该 seed 的主问题调度结果，可以得到单臂任务对应的可压缩时间 $T_S^{(r)}$、双臂协同任务对应的可压缩时间 $T_D^{(r)}$。因此，该 seed 下速度调整后的最大完工时间可表示为
\[
C_{\max}^{(r)}(s_S,s_D)
=
\frac{T_S^{(r)}}{s_S}
+
\frac{T_D^{(r)}}{s_D}.
\]
其中，$T_S^{(r)}$ 和 $T_D^{(r)}$ 随 seed 改变而改变，反映不同物块空间分布对任务时间结构的影响。

类似地，设 seed $r$ 下单臂任务基准能耗为 $E_S^{0,(r)}$，双臂协同任务基准能耗为 $E_D^{0,(r)}$。速度调整后的能耗函数可写为
\[
E^{(r)}(s_S,s_D)
=
E_S^{0,(r)}
\left[
\frac{\lambda_S}{s_S}
+
(1-\lambda_S)
\left((1-\rho_S)+\rho_S s_S^2\right)
\right]
+
E_D^{0,(r)}
\left[
\frac{\lambda_D}{s_D}
+
(1-\lambda_D)
\left((1-\rho_D)+\rho_D s_D^2\right)
\right].
\]
其中，$\lambda_S,\lambda_D$ 用于表示低速时任务持续时间增加带来的保持能耗影响，$\rho_S,\rho_D$ 用于表示高速时驱动能耗增加的影响。不同 seed 的能耗差异通过 $E_S^{0,(r)}$ 和 $E_D^{0,(r)}$ 体现。

为了保证同一组速度对所有 seed 都有效，需要对每一个 seed 都施加截止时间约束：
\[
C_{\max}^{(r)}(s_S,s_D)\leq D,
\qquad
\forall r\in\mathcal{R}.
\]
该约束表示：无论物块随机分布如何变化，只要属于同一任务数量组合，统一速度方案都必须满足给定截止时间 $D$。

本问题的目标是最小化最坏情形能耗。直接写可以表示为
\[
\min_{s_S,s_D}
\quad
\max_{r\in\mathcal{R}} E^{(r)}(s_S,s_D).
\]
该目标函数表示：在所有 seed 中找到能耗最大的那一个，并使这个最坏能耗尽可能小。\textbf{代码引入辅助变量 $z$，把不便直接求解的 min--max 目标改写为共同能耗上界约束。}
\[
E^{(r)}(s_S,s_D)\leq z,
\qquad
\forall r\in\mathcal{R}.
\]
这样，最小化最坏情形能耗就等价于最小化 $z$。

因此，问题二的鲁棒速度优化模型可写为
\[
\min_{s_S,s_D,z}
\quad
z
\]
\[
\mathrm{s.t.}\quad
C_{\max}^{(r)}(s_S,s_D)\leq D,
\qquad
\forall r\in\mathcal{R},
\]
\[
E^{(r)}(s_S,s_D)\leq z,
\qquad
\forall r\in\mathcal{R},
\]
\[
s_{\min}\leq s_S\leq s_{\max},
\qquad
s_{\min}\leq s_D\leq s_{\max}.
\]

从函数性质上看，$1/s$ 在 $s>0$ 上为凸函数，因此每个 seed 的时间函数 $C_{\max}^{(r)}(s_S,s_D)$ 是凸函数；能耗函数由 $1/s$ 项和 $s^2$ 项加权构成，在参数非负时也具有凸性。最坏情形目标本质上是多个凸能耗函数的最大值，而凸函数的最大值仍为凸函数。引入变量 $z$ 后，原来的问题被转化为带一组能耗上界约束的标准形式，便于后续使用逐次线性规划法进行数值求解。


\subsubsection{实现过程}

鲁棒速度模式选择问题采用逐次线性规划法求解。与问题一不同，问题二不再只处理平均意义下的一组完工时间和能耗，而是同时考虑同一任务数量组合下的多个 seed。每个 seed 都会产生一组截止时间约束和能耗上界约束，因此该问题虽然决策变量仍然较少，但约束数量明显增加。

本问题的原始目标可以写成最小化所有 seed 中的最大能耗：
\[
\min_{s_S,s_D}
\quad
\max_{r\in\mathcal{R}} E^{(r)}(s_S,s_D).
\]
其中，$\mathcal{R}$ 表示 seed 集合，$E^{(r)}(s_S,s_D)$ 表示 seed $r$ 下的能耗函数。直接处理 $\max$ 函数会带来两个问题：一是目标函数在多个 seed 同时达到最大值时可能不可微；二是后续 KKT 验证和数值求解都不方便。因此，本文引入辅助变量 $z$，将 min--max 问题改写为上界形式：
\[
E^{(r)}(s_S,s_D)\leq z,
\qquad
\forall r\in\mathcal{R}.
\]
这样，最小化最坏情形能耗就等价于最小化 $z$。改写后的目标函数变为线性函数：
\[
\min_{s_S,s_D,z}\quad z.
\]
这种处理相当于用变量 $z$ 表示所有 seed 能耗的共同上界。只要 $z$ 越小，所有 seed 中的最大能耗就越小。

令
\[
x
=
\begin{bmatrix}
	s_S\\
	s_D\\
	z
\end{bmatrix}.
\]
则鲁棒问题中的约束可以统一写成
\[
g_q(x)\geq 0,
\qquad q=1,2,\ldots,Q.
\]
其中，$Q$ 表示所有约束的总数，包括所有 seed 的截止时间约束、所有 seed 的能耗上界约束，以及速度上下界约束。具体来说，对每个 seed $r$，有截止时间约束
\[
g_T^{(r)}(x)
=
D-C_{\max}^{(r)}(s_S,s_D)
\geq 0,
\]
以及能耗上界约束
\[
g_E^{(r)}(x)
=
z-E^{(r)}(s_S,s_D)
\geq 0.
\]
此外还有速度边界约束
\[
s_S-s_{\min}\geq 0,
\qquad
s_{\max}-s_S\geq 0,
\]
\[
s_D-s_{\min}\geq 0,
\qquad
s_{\max}-s_D\geq 0.
\]

从函数性质看，本问题适合使用逐次线性规划法。首先，$s_S$ 和 $s_D$ 均有正的下界，因此 $1/s_S$、$1/s_D$、$s_S^2$ 和 $s_D^2$ 在可行域内都是连续可导函数。其次，时间函数中含有 $1/s$ 项，能耗函数中含有 $1/s$ 和 $s^2$ 项，这些函数在正数范围内具有凸性。因此，原问题可以看作一个具有凸结构的非线性规划问题。再次，变量只有 $s_S,s_D,z$ 三个，但约束来自多个 seed，逐次线性规划法可以在保留多约束结构的同时，把每一步转化为小规模线性规划子问题，适合本文自编程序实现。

逐次线性规划法的核心思想是：在当前迭代点附近，用一阶 Taylor 展开近似原来的非线性约束，从而得到一个线性规划子问题。设第 $k$ 次迭代的当前点为
\[
x^{(k)}
=
\begin{bmatrix}
	s_S^{(k)}\\
	s_D^{(k)}\\
	z^{(k)}
\end{bmatrix}.
\]
对于任意一个非线性约束 $g_q(x)\geq 0$，在 $x^{(k)}$ 处进行一阶 Taylor 展开：
\[
g_q(x)
\approx
g_q(x^{(k)})
+
\nabla g_q(x^{(k)})^T
\left(
x-x^{(k)}
\right).
\]
于是，在第 $k$ 次迭代中，用线性约束
\[
g_q(x^{(k)})
+
\nabla g_q(x^{(k)})^T
\left(
x-x^{(k)}
\right)
\geq 0
\]
近似原来的非线性约束。

需要注意的是，线性化只在当前点附近可靠。\textbf{为避免候选解在线性近似中可行、代回原模型却不可行，代码在逐次线性规划中加入步长限制，即信赖区域约束：}
\[
|s_S-s_S^{(k)}|\leq \Delta_k,
\]
\[
|s_D-s_D^{(k)}|\leq \Delta_k.
\]
其中，$\Delta_k$ 表示第 $k$ 次迭代的信赖区域大小。该限制保证每一次线性近似只在当前点附近使用，从而提高迭代稳定性。

因此，第 $k$ 次迭代需要求解的线性规划子问题为
\[
\begin{aligned}
	\min_{s_S,s_D,z}\quad
	& z\\
	\mathrm{s.t.}\quad
	& g_q(x^{(k)})
	+
	\nabla g_q(x^{(k)})^T
	\left(
	x-x^{(k)}
	\right)
	\geq 0,
	\qquad q=1,2,\ldots,Q,\\
	& |s_S-s_S^{(k)}|\leq \Delta_k,\\
	& |s_D-s_D^{(k)}|\leq \Delta_k,\\
	& s_{\min}\leq s_S\leq s_{\max},\\
	& s_{\min}\leq s_D\leq s_{\max}.
\end{aligned}
\]
该子问题的目标函数是线性的，所有约束也都是线性的，因此是一个标准线性规划问题。

\textbf{线性规划子问题求解后，程序不会直接接受候选解，而是将其代回原非线性模型进行可行性复核。}若候选解满足所有 seed 的原始截止时间约束：
\[
C_{\max}^{(r)}(s_S,s_D)\leq D,
\qquad
\forall r\in\mathcal{R},
\]
同时满足所有 seed 的能耗上界约束：
\[
E^{(r)}(s_S,s_D)\leq z,
\qquad
\forall r\in\mathcal{R},
\]
并且目标值 $z$ 相比当前解有所改善，则接受该候选解作为新的迭代点。\textbf{若回代后不可行或目标改善不明显，程序便缩小信赖区域 $\Delta_k$ 并重新求解，从而形成“线性化--回代验证--自适应缩步”的闭环。}

\textbf{利用线性规划子问题只有 $s_S,s_D,z$ 三个变量的特点，代码采用顶点枚举自行求解：枚举约束交点、筛选可行顶点，再选取 $z$ 最小者。}这样可以避免依赖外部大型优化求解器，也与本文主问题中自编算法的思路保持一致。

这个过程如图 \ref{fig:robust-slp-flow} 所示。

\begin{figure}[H]
	\centering
	\begin{tikzpicture}[
		node distance=0.42cm,
		flowbox/.style={
			rectangle,
			rounded corners=2pt,
			draw=gray!60,
			line width=0.4pt,
			minimum height=0.72cm,
			text width=13.6cm,
			align=left,
			inner xsep=8pt,
			inner ysep=5pt,
			font=\small
		},
		arrow/.style={
			-{Stealth[length=2mm,width=1.5mm]},
			draw=gray!70,
			line width=0.4pt
		}
		]
		
		\node[flowbox, fill=purple!8] (step1)
		{\textbf{读取多 seed 数据}\quad 读取同一 \texttt{counts\_code} 下所有 seed 的主问题结果，分别构造 $C_{\max}^{(r)}(s_S,s_D)$ 和 $E^{(r)}(s_S,s_D)$};
		
		\node[flowbox, fill=green!8, below=of step1] (step2)
		{\textbf{引入辅助变量}\quad 引入辅助变量 $z$，将最小化最坏情形能耗问题转化为 $\min z$ 及一组能耗上界约束};
		
		\node[flowbox, fill=purple!8, below=of step2] (step3)
		{\textbf{选取初始点}\quad 选取初始速度 $s_S^{(0)},s_D^{(0)}$，并令 $z^{(0)}$ 为该速度下所有 seed 能耗的最大值};
		
		\node[flowbox, fill=green!8, below=of step3] (step4)
		{\textbf{约束线性化}\quad 在当前点 $x^{(k)}$ 处，对所有非线性约束 $g_q(x)\geq 0$ 作一阶 Taylor 线性化};
		
		\node[flowbox, fill=purple!8, below=of step4] (step5)
		{\textbf{加入信赖区域}\quad 加入信赖区域约束 $|s_S-s_S^{(k)}|\leq \Delta_k$ 和 $|s_D-s_D^{(k)}|\leq \Delta_k$};
		
		\node[flowbox, fill=green!8, below=of step5] (step6)
		{\textbf{求解线性规划子问题}\quad 求解当前线性规划子问题，得到候选解 $\hat{x}^{(k)}$};
		
		\node[flowbox, fill=purple!8, below=of step6] (step7)
		{\textbf{原模型可行性检查}\quad 将候选解代回原非线性约束中检查可行性，并判断最坏情形能耗是否改善};
		
		\node[flowbox, fill=green!8, below=of step7] (step8)
		{\textbf{接受或调整步长}\quad 若候选解可行且改善有效，则接受该解，并可适当扩大或保持信赖区域；若候选解不可行或改善不足，则缩小信赖区域并重新求解};
		
		\node[flowbox, fill=purple!8, below=of step8] (step9)
		{\textbf{迭代直到收敛}\quad 重复上述过程，直到 $z$ 的变化量、速度变量变化量或信赖区域大小满足停止条件};
		
		\draw[arrow] (step1) -- (step2);
		\draw[arrow] (step2) -- (step3);
		\draw[arrow] (step3) -- (step4);
		\draw[arrow] (step4) -- (step5);
		\draw[arrow] (step5) -- (step6);
		\draw[arrow] (step6) -- (step7);
		\draw[arrow] (step7) -- (step8);
		\draw[arrow] (step8) -- (step9);
		
	\end{tikzpicture}
	\caption{逐次线性规划法求解鲁棒速度模式选择问题的流程}
	\label{fig:robust-slp-flow}
\end{figure}

\textbf{该方法没有把多个 seed 简单平均，而是完整保留每个 seed 的时间约束和能耗约束，再通过逐次线性化将鲁棒非线性问题转化为一系列小规模线性规划。}最终程序输出统一速度倍率 $s_S^{*},s_D^{*}$、最坏情形能耗 $z^{*}$、各 seed 的可行性检查结果以及 KKT 验证结果。由此可以进一步比较二臂和三臂模式在随机任务分布扰动下的稳健性差异。


\subsubsection{KKT验证条件}

鲁棒速度模式选择问题的 KKT 条件比平均速度--能耗优化问题更加复杂。其标准形式为
\[
\begin{aligned}
    \min_x\quad & z\\
    \text{s.t.}\quad & g_q(x)\geq 0,\qquad q=1,2,\ldots,Q.
\end{aligned}
\]
其中，$x=(s_S,s_D,z)^T$，$Q$ 包含所有 seed 的能耗上界约束、截止时间约束以及速度上下界约束。

其 KKT 条件可以写为
\[
\nabla z-\sum_{q=1}^{Q}\mu_q\nabla g_q(x^*)=0,
\]
\[
g_q(x^*)\geq 0,\qquad q=1,2,\ldots,Q,
\]
\[
\mu_q\geq 0,\qquad q=1,2,\ldots,Q,
\]
\[
\mu_q g_q(x^*)=0,\qquad q=1,2,\ldots,Q.
\]
其中，平稳性条件要求目标变量 $z$ 与各个 seed 的能耗约束、时间约束以及速度边界约束之间形成一阶平衡关系；互补松弛条件用于判断哪些 seed 或边界约束在最优点附近真正起作用。

\textbf{程序先用原非线性约束保证鲁棒可行性，再用 KKT 检查判断是否满足严格的一阶局部最优性。}若部分行未通过 KKT 检查，则应将其解释为鲁棒可行解或数值近似解，而不能直接称为严格最优解。
